% Minimal preamble for Open Logic Project sources (pinned commit 1e960beff9ed7835bf3e3f1335e21af3439cd107).
%
% OLP's own master preamble cannot run on the project host: TeX Live 2025 there lacks
% prooftree.sty, ntheorem.sty, extarrows.sty, mathdots.sty, isomath.sty and subfiles.sty.
% Following the project's own rule ("if the author preamble needs packages the toolchain lacks,
% do not invent macros -- use a minimal working preamble"), this file declares, verbatim in
% spirit and numbering behaviour, exactly the theorem-like environments OLP defines in
% sty/open-logic-book-envs.sty (thm/lem/prop/cor/defn/ex/prob/rem/note/case/axiom/conv),
% plus the shorthand accents OLP uses (`!!a`, `!!b`, `!!{...}`).
%
% Author text is untouched; only the packaging is ours.
\documentclass[10pt,a4paper]{article}

\usepackage{amsmath}
\usepackage{amssymb}
\usepackage{amsthm}
\usepackage{graphicx}
\usepackage{hyperref}
\usepackage{enumitem}
\usepackage[framemethod=default]{mdframed}

% ---- OLP shorthand accents (from sty/open-logic.sty) ----
% NB: \! is taken by LaTeX (negative thin space), so we only define the !!a / !!b families
% and the \olup helper used by them.
\newcommand*{\olup}[1]{\MakeUppercase{#1}}
\def\!!!a#1{\olup{#1}}
\def\!!a#1{\olup{#1}}
\def\!!b#1{\olup{#1}}

% ---- theorem-like environments, same family numbering as OLP ----
\theoremstyle{plain}
\newtheorem{thm}{Theorem}[section]
\newtheorem{lem}[thm]{Lemma}
\newtheorem{prop}[thm]{Proposition}
\newtheorem{cor}[thm]{Corollary}
\newtheorem{axiom}{Axiom}

\theoremstyle{definition}
\newtheorem{defn}[thm]{Definition}
\newtheorem{ex}[thm]{Example}
\newtheorem{prob}{Problem}[section]
\newtheorem{probd}{Problem}
\newtheorem{conv}[thm]{Convention}

\theoremstyle{remark}
\newtheorem{rem}[thm]{Remark}
\newtheorem{note}[thm]{Note}
\newtheorem{case}[thm]{Case}

% ---- environments OLP defines as wrappers ----
\newenvironment{probtag}[1]{\begin{prob}}{\end{prob}}
\newenvironment{explain}{}{}
\newenvironment{explain*}{}{}
\newenvironment{editorial}{}{}
\newenvironment{digress}{}{}
\newenvironment{digress*}{}{}
\newenvironment{history}{}{}
\newenvironment{reference}{}{}
\newcommand*{\olfileid}[3]{}
\newcommand*{\olsection}[1]{\section*{#1}}
\newcommand*{\olchapter}[3]{}
\newcommand*{\olpart}[2]{}
\newcommand*{\olimport}[1]{}
\newcommand*{\olinclude}[1]{}
\newcommand*{\olpath}[1]{}
\newcommand*{\ollangid}[1]{}
\newcommand*{\problemtag}[1]{}
\newcommand*{\problemauthor}[1]{}
\newcommand*{\ollabel}[1]{\label{#1}}
% \olref takes 0-3 OPTIONAL arguments plus one mandatory label:
%     \olref{lab}  |  \olref[from]{lab}  |  \olref[from][to]{lab}  |  \olref[a][b][c]{lab}
% A single-argument definition silently swallows "[sth]" and emits a broken reference
% (measured on the host: the text came out as "??iso]x"). The chain below consumes the
% brackets and cross-references only the mandatory argument.
\makeatletter
\newcommand{\olref}{\@ifnextchar[{\olref@i}{\olref@final}}
\def\olref@i[#1]{\@ifnextchar[{\olref@ii}{\olref@final}}
\def\olref@ii[#1]{\@ifnextchar[{\olref@iii}{\olref@final}}
\def\olref@iii[#1]{\olref@final}
\def\olref@final#1{\ref{#1}}
\makeatother
\newcommand*{\olnewpage}{\clearpage}

% ===================== OLP compatibility layer (1) =====================

\makeatletter

% ---------------- proof trees (OLP's prooftree is a bussproofs derivative) ----------------
\providecommand{\AxiomC}[1]{}
\providecommand{\UnaryInfC}[1]{}
\providecommand{\BinaryInfC}[1]{}
\providecommand{\TrinaryInfC}[1]{}
\providecommand{\QuaternaryInfC}[1]{}
\providecommand{\QuinaryInfC}[1]{}
\providecommand{\RightLabel}[1]{}
\providecommand{\LeftLabel}[1]{}
\providecommand{\doubleLine}{}
\providecommand{\noLine}{}
\providecommand{\DisplayProof}{}
\providecommand{\fCenter}{}
\newenvironment{prooftree}{}{}
\newenvironment{ltprooftree}{}{}
\newenvironment{oltableau}{}{}

% ---------------- logical notation ----------------
\providecommand{\lif}{\mathbin{\rightarrow}}
\providecommand{\lforall}{\mathord{\forall}}
\providecommand{\lexists}{\mathord{\exists}}
\providecommand{\ltrue}{\top}
\providecommand{\lfalse}{\bot}
\providecommand{\Ltrue}{\top}
\providecommand{\Lfalse}{\bot}
\providecommand{\Sequent}{\mathrel{\Rightarrow}}
\providecommand{\Proves}{\mathrel{\vdash}}
\providecommand{\Entails}{\mathrel{\models}}
\providecommand{\Subst}[3]{#1[#2/#3]}
\providecommand{\sFmla}[2]{#1(#2)}
\providecommand{\pFmla}[2]{#1(#2)}
\providecommand{\formula}{\relax}
\providecommand{\True}{\top}
\providecommand{\False}{\bot}
\providecommand{\Args}[1]{#1}
\providecommand{\Ts}{\ensuremath{\top}}
\providecommand{\Fi}{\ensuremath{\bot}}
\providecommand{\lfalsum}{\bot}

% ---------------- objects, structures, metatheory ----------------
\providecommand{\Obj}{}
\providecommand{\Struct}{\ensuremath{\mathfrak{A}}}
\providecommand{\Domain}{\ensuremath{|\mathfrak{A}|}}
\providecommand{\Assign}{\ensuremath{s}}
\providecommand{\Sat}[2]{\ensuremath{#1 \models #2}}
\providecommand{\mSat}[3]{\ensuremath{#1 \models #2[#3]}}
\providecommand{\mModel}[2]{\ensuremath{#1 \models #2}}
\providecommand{\Lang}{\ensuremath{\mathcal{L}}}
\providecommand{\Th}{\ensuremath{\mathrm{Th}}}
\providecommand{\Log}{\ensuremath{\mathrm{L}}}
\providecommand{\tuple}[1]{\langle #1 \rangle}
\providecommand{\num}[1]{\underline{#1}}
\providecommand{\cfind}[1]{\varphi_{#1}}
\providecommand{\Nat}{\ensuremath{\mathbb{N}}}
\providecommand{\Mult}{\ensuremath{\mathrm{Mult}}}
\providecommand{\Add}{\ensuremath{\mathrm{Add}}}
\providecommand{\Exp}{\ensuremath{\mathrm{Exp}}}
\providecommand{\Zero}{\ensuremath{0}}
\providecommand{\Succ}{\ensuremath{\mathrm{S}}}
\providecommand{\lambd}{\ensuremath{\lambda}}
\providecommand{\bredpar}{\mathrel{\rightarrow_\beta}}
\providecommand{\bcd}[1]{\ensuremath{\ulcorner #1 \urcorner}}
\providecommand{\Set}[1]{\ensuremath{\{#1\}}}
\providecommand{\indcase}[1]{\par\noindent\emph{Case #1}.\quad}
\providecommand{\indfrm}[1]{\par\noindent\emph{Form #1}.\quad}
\providecommand{\iftag}[1]{#1}
\providecommand{\tagitem}[1]{\item[#1]}
\providecommand{\Elim}[1]{}
\providecommand{\Intro}[1]{}
\providecommand{\LeftR}[1]{}
\providecommand{\RightR}[1]{}
\providecommand{\TAss}[1]{}
\providecommand{\TRule}[1]{}

\makeatother
\providecommand{\olpath}{.}

% ===================== OLP compatibility layer (2) =====================
% --- second compatibility block: notation and shorthands used by the frozen OLP units ---
\providecommand{\CutCS}{\mathrm{Cut}}
\providecommand{\LeftSubproofLabel}{\relax}
\providecommand{\RightSubproofLabel}{\relax}
\providecommand{\maeh}{\relax}
\providecommand{\cheight}{\relax}
\providecommand{\ande}{\mathrm{and}_E}
\providecommand{\andi}{\mathrm{and}_I}
\providecommand{\ore}{\mathrm{or}_E}
\providecommand{\ori}{\mathrm{or}_I}
\providecommand{\typeof}{\mathord{:}}
\providecommand{\pValue}[1]{#1}
\providecommand{\Sent}{\mathrm{Sent}}
\providecommand{\cardplus}{\mathbin{+}}
\providecommand{\cardtimes}{\mathbin{\cdot}}
\providecommand{\cardexpo}{\mathbin{\hat{\ }}}
\providecommand{\cardnless}{\not<}
\providecommand{\cardle}{\leq}
\providecommand{\funfromto}[3]{#1\colon #2 \to #3}
\providecommand{\funimage}[1]{\operatorname{im}(#1)}
\providecommand{\fdefined}{\mathord{\downarrow}}
\providecommand{\canonord}{\operatorname{ord}}
\providecommand{\ordtimes}{\mathbin{\cdot}}
\providecommand{\ordeq}[2]{#1 = #2}
\providecommand{\ordneq}[2]{#1 \neq #2}
\providecommand{\ordtype}[1]{\operatorname{ord}(#1)}
\providecommand{\ordsucc}{\mathord{+1}}
\providecommand{\ordunion}{\cup}
\providecommand{\onesphere}{S^1}
\providecommand{\rotationsgroup}{\mathrm{SO}(2)}
\providecommand{\equivrep}[1]{[#1]}
\providecommand{\rlexless}{\prec}
\providecommand{\ZFC}{\mathrm{ZFC}}
\providecommand{\ZF}{\mathrm{ZF}}
\providecommand{\Rat}{\mathbb{Q}}
\providecommand{\Real}{\mathbb{R}}
\providecommand{\PosInt}{\mathbb{Z}^{+}}
\providecommand{\Id}{\mathrm{id}}
\providecommand{\Bin}{\mathrm{Bin}}
\providecommand{\pto}{\rightharpoonup}
\providecommand{\dom}{\operatorname{dom}}
\providecommand{\ran}{\operatorname{ran}}
\providecommand{\len}{\operatorname{len}}
\providecommand{\concat}{\mathbin{{}^\frown}}
\providecommand{\card}[1]{\lvert #1 \rvert}
\providecommand{\Pow}{\mathcal{P}}
\providecommand{\Setabs}[2]{\{\, #1 \mid #2 \,\}}
\providecommand{\comp}{\circ}
\providecommand{\Char}{\chi}
\providecommand{\Prf}{\operatorname{Prf}}
\providecommand{\Prov}{\operatorname{Prov}}
\providecommand{\Proj}{\operatorname{proj}}
\providecommand{\Gn}{\operatorname{gn}}
\providecommand{\gn}{\operatorname{gn}}
\providecommand{\Value}[1]{\operatorname{val}(#1)}
\providecommand{\Deduce}{\mathrm{D}}
\providecommand{\Dual}{\mathrm{d}}
\providecommand{\Ax}{\mathrm{Ax}}
\providecommand{\pSat}[3]{#1 \models #2\,[#3]}
\providecommand{\pAssign}{\sigma}
\providecommand{\bforall}{\forall}
\providecommand{\bexists}{\exists}
\providecommand{\liff}{\mathbin{\leftrightarrow}}
\providecommand{\FV}{\operatorname{FV}}
\providecommand{\UnaryInf}{\ensuremath{\mid}}
\providecommand{\BinaryInf}{\ensuremath{\mid}}
\providecommand{\depth}{\operatorname{depth}}
\providecommand{\oliflabeldef}[3]{#1}
\providecommand{\TMblank}{\sqcup}
\providecommand{\TMstroke}{\mid}
\providecommand{\TMstay}{\mathrm{S}}
\providecommand{\TMright}{\mathrm{R}}
\providecommand{\TMleft}{\mathrm{L}}
\providecommand{\TMendtape}{\triangleright}
\providecommand{\umin}{\mu}
\providecommand{\Sum}{\sum}
\providecommand{\citep}[1]{\cite{#1}}
\providecommand{\citet}[1]{\cite{#1}}
\providecommand{\citeauthor}[1]{\cite{#1}}
\providecommand{\citeyear}[1]{\cite{#1}}

% ===================== OLP compatibility layer (3) =====================
% --- third compatibility block: remaining notation seen in the frozen OLP units ---
\providecommand{\iso}{\cong}
\providecommand{\nszero}{\mathrm{0}}
\providecommand{\nssucc}{\mathrm{S}}
\providecommand{\nsplus}{\mathbin{+}}
\providecommand{\nssless}{<}
\providecommand{\nsless}{<}
\providecommand{\PL}{\mathrm{PL}}
\providecommand{\Nec}{\mathord{\Box}}
\providecommand{\RK}{\mathrm{RK}}
\providecommand{\ST}{\mathrm{ST}}
\providecommand{\SSubst}[3]{#1[#2/#3]}
\providecommand{\subst}[3]{#1[#2/#3]}
\providecommand{\tagfalse}{\mathrm{F}}
\providecommand{\cutr}{\mathrm{Cut}}
\providecommand{\cardneq}{\neq}
\providecommand{\issucc}[1]{\operatorname{succ}(#1)}
\providecommand{\tup}[1]{\langle #1 \rangle}
\providecommand{\Var}{\mathrm{Var}}
\providecommand{\Bew}{\operatorname{Bew}}
\providecommand{\Con}{\operatorname{Con}}
\providecommand{\Cons}{\operatorname{Cons}}
\providecommand{\pDeduce}{\mathrm{D}}
\providecommand{\psat}{\models}
\providecommand{\den}[1]{\llbracket #1 \rrbracket}
\providecommand{\Sem}[2]{\llbracket #1 \rrbracket_{#2}}
\providecommand{\DomainOf}[1]{\lvert #1 \rvert}
\providecommand{\Form}{\mathrm{Form}}
\providecommand{\Terms}{\mathrm{Term}}
\providecommand{\Atom}{\mathrm{Atom}}
\providecommand{\FVars}[1]{\operatorname{FV}(#1)}
\providecommand{\fv}[1]{\operatorname{fv}(#1)}
\providecommand{\tuplep}[1]{\langle #1 \rangle}
\providecommand{\restrict}[2]{#1 \upharpoonright #2}
\providecommand{\sto}{\rightarrow}
\providecommand{\FmlaSet}[1]{\mathrm{Fmla}_{#1}}
\providecommand{\Thm}{\mathrm{Thm}}
\providecommand{\Eval}{\operatorname{val}}
\providecommand{\blank}{\sqcup}
\providecommand{\blankTape}{\sqcup}
\providecommand{\Consis}{\operatorname{Con}}
\providecommand{\drs}{\mathrm{d}}
\providecommand{\VRule}{\mid}
\providecommand{\Bet}{\mathrm{B}}
\providecommand{\Seg}[1]{\mathrm{Seg}_{#1}}
\providecommand{\Enums}{\mathrm{En}}
\providecommand{\Z}{\mathbb{Z}}
\providecommand{\N}{\mathbb{N}}
\providecommand{\abs}[1]{\lvert #1 \rvert}
\providecommand{\floor}[1]{\lfloor #1 \rfloor}
\providecommand{\ceil}[1]{\lceil #1 \rceil}

% ===================== OLP compatibility layer (4) =====================
% --- fourth compatibility block: harvested from the per-candidate compile logs ---
\providecommand{\beredpar}{\mathrel{\rightarrow_\beta}}
\makeatletter
\def\bered{\@ifnextchar[{\bered@i}{\mathrel{\rightarrow_\beta}}}
\def\bered@i[#1]{\mathrel{\rightarrow_\beta}}
\def\bred{\@ifnextchar[{\bred@i}{\mathrel{\rightarrow_\beta}}}
\def\bred@i[#1]{\mathrel{\rightarrow_\beta}}
\makeatother
\providecommand{\beredone}{\mathrel{=_\beta}}
\providecommand{\bredone}{\mathrel{=_\beta}}
\providecommand{\xredone}{\mathrel{=_\beta}}
\providecommand{\aconvone}{\mathrel{=_\alpha}}
\providecommand{\cardless}{<}
\providecommand{\tagrefs}[1]{#1}
\providecommand{\ordplus}{\mathbin{+}}
\providecommand{\funrestrictionto}[2]{#1 \upharpoonright #2}
\providecommand{\fundefined}{\mathord{\uparrow}}
\providecommand{\Trm}{\mathrm{Trm}}
\providecommand{\OProv}{\operatorname{Prov}}
\providecommand{\OPrf}{\operatorname{Prf}}
\providecommand{\aeq}{\mathrel{=}}
\providecommand{\mClass}[1]{[#1]}
\providecommand{\equivclass}[1]{[#1]}
\providecommand{\pheight}{\operatorname{ht}}
\providecommand{\phy}{\operatorname{ht}}
\providecommand{\Cut}{\mathrm{Cut}}
\providecommand{\QuantRank}{\operatorname{qr}}
\providecommand{\elemequiv}{\mathrel{\equiv}}
\providecommand{\PIso}{\mathrm{PIso}}
\providecommand{\Theory}{\mathrm{Th}}
\providecommand{\supstrict}{\subsetneq}
\providecommand{\trcl}{\operatorname{trcl}}
\providecommand{\cref}[1]{\ref{#1}}
\providecommand{\Discharge}{\mathrm{Dis}}
\providecommand{\Weakening}{\mathrm{Wk}}
\providecommand{\ZFminus}{\mathrm{ZF}^{-}}
\providecommand{\RK}{\mathrm{RK}}
\makeatletter
\renewcommand{\Dual}{\@ifnextchar[{\Dual@i}{\ifmmode\mathrm{d}\else$\mathrm{d}$\fi}}
\def\Dual@i[#1]{\ifmmode\mathrm{d}[#1]\else$\mathrm{d}[#1]$\fi}
\makeatother
\providecommand{\RightR}{\mathrm{R}}
\providecommand{\LeftR}{\mathrm{L}}
\providecommand{\Intro}{\mathrm{I}}
\providecommand{\Axiom}{\mathrm{Ax}}
\providecommand{\formula}[1]{#1}
\providecommand{\stared}{\ast}
\providecommand{\oset}[2]{\{\, #1 \mid #2 \,\}}
\providecommand{\setof}[2]{\{\, #1 \mid #2 \,\}}
\providecommand{\pow}[1]{\mathcal{P}(#1)}
\providecommand{\ordpair}[2]{\langle #1, #2 \rangle}
\providecommand{\ordtriple}[3]{\langle #1, #2, #3 \rangle}
\providecommand{\union}{\cup}
\providecommand{\inter}{\cap}
\providecommand{\seq}[1]{\langle #1 \rangle}
\providecommand{\SuchThat}{\mid}
\providecommand{\where}{\mid}
\providecommand{\TVar}{\mathrm{TVar}}
\providecommand{\Bool}{\mathrm{Bool}}
\providecommand{\Val}{\operatorname{Val}}
\providecommand{\FV}{\operatorname{FV}}
\providecommand{\TL}[1]{\mathrm{TL}}
\providecommand{\xstared}{\ast}

% ===================== OLP compatibility layer (5) =====================
% --- fifth compatibility block: rescues the quarantined OLP candidates ---
\providecommand{\ident}{\operatorname{id}}
\providecommand{\setrank}{\operatorname{rank}}
\providecommand{\cardfont}{\mathfrak{c}}
\providecommand{\Complement}{\overline}
\providecommand{\QR}{\mathrm{QR}}
\providecommand{\becd}{\ensuremath{\ulcorner}}
\providecommand{\disjointsum}{\mathbin{\oplus}}
\providecommand{\cardeq}[2]{#1 = #2}
\providecommand{\xred}{\mathrel{\rightarrow}}
\providecommand{\cardless}{<}
\providecommand{\cardplus}{\mathbin{+}}
\providecommand{\cardtimes}{\mathbin{\cdot}}
\providecommand{\cardexpo}{\mathbin{\hat{\ }}}
\providecommand{\cardneq}{\neq}
\providecommand{\cardle}{\leq}
\providecommand{\cardnless}{\not<}
% --- sixth compatibility block: names reported by real compile logs ---
\providecommand{\ORefut}{\operatorname{Refut}}
\providecommand{\ORProv}{\operatorname{Prov}}
\providecommand{\cutrank}{\operatorname{cr}}
\providecommand{\Refut}{\operatorname{Refut}}
\providecommand{\DischargeRule}{\mathrm{Dis}}
\providecommand{\Contraction}{\mathrm{C}}
\providecommand{\bottomAlignProof}{\relax}
\providecommand{\FalseCl}{\bot}
\providecommand{\red}{\mathrel{\rightarrow}}
\providecommand{\varAssign}{\sigma}
\providecommand{\defiff}{\mathrel{\leftrightarrow}}
\providecommand{\FalseInt}{\bot}
\providecommand{\fn}{\operatorname{fn}}
\providecommand{\stared}{\ast}
\providecommand{\dotted}{\cdot}
\providecommand{\Gamma}{\Gamma}
% --- sixth compatibility block: names reported by real compile logs ---
\providecommand{\ORefut}{\operatorname{Refut}}
\providecommand{\ORProv}{\operatorname{Prov}}
\providecommand{\cutrank}{\operatorname{cr}}
\providecommand{\Refut}{\operatorname{Refut}}
\providecommand{\DischargeRule}{\mathrm{Dis}}
\providecommand{\Contraction}{\mathrm{C}}
\providecommand{\bottomAlignProof}{\relax}
\providecommand{\FalseCl}{\bot}
\providecommand{\red}{\mathrel{\rightarrow}}
\providecommand{\varAssign}{\sigma}
\providecommand{\defiff}{\mathrel{\leftrightarrow}}
\providecommand{\FalseInt}{\bot}
\providecommand{\fn}{\operatorname{fn}}
\providecommand{\stared}{\ast}
\providecommand{\dotted}{\cdot}
\providecommand{\Gamma}{\Gamma}
\begin{document}

% ==== OLP file: content/lambda-calculus/church-rosser/beta-eta-reduction.tex ====
% Part: lambda-calculus
% Chapter: church-rosser
% Section: beta-eta-reduction



\olfileid{lam}{cr}{be}

\olsection{$\beta\eta$-reduction}

The Church--Rosser property holds for
$\beta\eta$-reduction ($\bered$).

\begin{lem}\ollabel{lem:one-par}
  If $M \beredone M'$, then $M \beredpar M'$.
\end{lem}

\begin{proof} 
  By induction on the !!{derivation} of $M \beredone M'$. If $M \bredone
  M'$ by $\eta$-conversion (i.e., \olref[syn][eta]{defn:beredone}), we
  use \olref[pbe]{thm:refl}. The other cases are as in
  \olref[b]{lem:one-par}.
\end{proof}


\begin{lem}\ollabel{lem:par-red}
  If $M \beredpar M'$, then $M \bered M'$.
\end{lem}

\begin{proof} Induction on the !!{derivation} of $M \beredpar M'$.

  If the last rule is \olref[pbe]{defn:beredpar5}, then $M$ is
  $\lambd[x][Nx]$ and $M'$ is $N'$ for some $x$, $N$, $N'$ where $x
  \notin FV(N)$ and $N \beredpar N'$. Thus we can first reduce
  $\lambd[x][Nx]$ to $N$ by $\eta$-conversion, followed
  by the series of $\beredone$ steps that show that $N \bered N'$,
  which  holds by induction hypothesis.
\end{proof}


\begin{lem}\ollabel{lem:str}
  $\bered$ is the smallest transitive relation containing $\beredpar$.
\end{lem}

\begin{proof}
  As in \olref[b]{lem:str}
\end{proof}

\begin{thm}\ollabel{thm:cr}
  $\bered$ satisfies Church--Rosser property.
\end{thm}

\begin{proof}
  By \olref[dap]{thm:str}, \olref[pbe]{thm:cr} and \olref{lem:str}.
\end{proof}
% ==== OLP file: content/lambda-calculus/church-rosser/beta-reduction.tex ====
% Part: lambda-calculus
% Chapter: church-rosser
% Section: beta-reduction



\olfileid{lam}{cr}{b}

\olsection{$\beta$-reduction}

\begin{lem}\ollabel{lem:one-par}
  If $M \bredone M'$, then $M \bredpar M'$.
\end{lem}
\begin{proof} If $M \bredone M'$, then $M$ is
    $(\lambd[x][N])Q$, $M'$ is $\Subst{N}{Q}{x}$, for some
    $x$, $N$, and~$Q$. Since $N \bredpar N$ and $Q \bredpar Q$ by
    \olref[pb]{thm:refl}, we immediately have $(\lambd[x][N])Q
    \bredpar \Subst{N}{Q}{x}$ by \olref[pb]{defn:bredpar}\olref[pb]{defn:bredpar4}.
\end{proof}

\begin{lem}\ollabel{lem:par-red}
  If $M \bredpar M'$, then $M \bred M'$.
\end{lem}

\begin{proof} By induction on the !!{derivation} of $M \bredpar M'$.
  \begin{enumerate}
  \item The last rule is \olref[pb]{defn:bredpar1}: Then $M$ and $M'$
    are just $x$, and $x \bred x$.
  \item The last rule is \olref[pb]{defn:bredpar2}: $M$ is 
    $\lambd[x][N]$ and $M'$ is $\lambd[x][N']$ for some $x$, $N$, $N'$, where
    $N \bredpar N'$. By induction hypothesis we have $N \bred N'$. Then
    $\lambd[x][N] \bred \lambd[x][N']$ (by the same series of
    $\bredone$ contractions as $N \bred N'$).
  \item The last rule is \olref[pb]{defn:bredpar3}: $M$ is 
    $PQ$ and $M'$ is $P'Q'$ for some $P$, $Q$, $P'$, $Q'$, where $P \bredpar P'$
    and $Q \bredpar Q'$. By induction hypothesis we have $P \bred P'$ and $Q \bred
    Q'$.  So $PQ \bred P'Q'$ by the reduction sequence $P \bred P'$ followed
    by the reduction $Q \bred Q'$.
  \item The last rule is \olref[pb]{defn:bredpar4}: $M$ is
    $(\lambd[x][N])Q$ and $M'$ is  $\Subst{N'}{Q'}{x}$ for some $x$,
    $N$, $M'$, $Q$, $Q'$, where $N \bredpar N'$ and $Q \bredpar Q'$.
    By induction hypothesis we get $Q \bred Q'$ and $N \bred N'$. So
    $(\lambd[x][N])Q \bred \Subst{N'}{Q'}{x}$ by $N \bred N'$ followed
    by $Q \bred Q'$ and finally contraction of $(\lambd[x][N'])Q'$ to
    $\Subst{N'}{Q'}{x}$.
  \end{enumerate}
\end{proof}

\begin{lem}\ollabel{lem:str}
  $\bred$ is the smallest transitive relation containing $\bredpar$.
\end{lem}

\begin{proof}
  Let $\xred$ be the smallest transitive relation containing
  $\bredpar$.

  $\bred \subseteq \xred$: Suppose $M \bred M'$, i.e., $M \ident M_1
  \bredone \dots \bredone M_k \ident M'$. By \olref{lem:one-par}, $M
  \ident M_1 \bredpar \dots \bredpar M_k \ident M'$. Since is $\xred$
  contains $\bredpar$ and is transitive, $M \xred M'$.

  $\xred \subseteq \bred$: Suppose $M \xred M'$, i.e., $M
  \ident M_1 \bredpar \dots \bredpar M_k \ident M'$. By \olref{lem:par-red}, $M \ident M_1
  \bred \dots \bred M_k \ident M'$. Since $\bred$ is transitive, $M
  \bred M'$.
\end{proof}

\begin{thm}\ollabel{thm:cr}
  $\bred$ satisfies the Church--Rosser property.
\end{thm}

\begin{proof}
  Immediate from \olref[dap]{thm:str}, \olref[pb]{thm:cr}, and \olref{lem:str}.
\end{proof}


% ==== OLP file: content/lambda-calculus/church-rosser/church-rosser.tex ====
% Part: lambda-calculus
% Chapter: church-rosser



\olchapter{lam}{cr}{The Church--Rosser Property}

\olimport{definitions-and-properties}
\olimport{parallel-beta-reduction}
\olimport{beta-reduction}
\olimport{parallel-beta-eta-reduction}
\olimport{beta-eta-reduction}

\OLEndChapterHook



% ==== OLP file: content/lambda-calculus/church-rosser/definitions-and-properties.tex ====
% Part: lambda-calculus
% Chapter: church-rosser
% Section: definitions-and-properties



\olfileid{lam}{cr}{dap}

\olsection{Definition and Properties}

In this chapter we introduce the concept of Church--Rosser property and
some common properties of this property.

\begin{defn}[Church--Rosser property, CR] A relation $\xredone$ on
  terms is said to satisfy the \emph{Church--Rosser property} iff,
  whenever $M \xredone P$ and $M \xredone Q$, then there exists
  some~$N$ such that $P \xredone N$ and $Q \xredone N$.
\end{defn}

We can view the lambda calculus as a model of computation in which
terms in normal form are ``values'' and a reducibility relation on
terms are the ``calculation rules.'' The Church--Rosser property states
is that when there is more than one way to proceed with a calculation,
there is still only a single value of the expression.

To take an example from elementary algebra, there's more than one way
to calculate $4 \times (1+2) + 3$. It can either be reduced to $4
\times 3+3$ (if we first reduce $1+2$ to~$3$) or to $4 \times 1+4
\times 2+3$ (if we first reduce $4 \times (1+2)$ using
distributivity). Both of these, however, can be further reduced to
$12+3$.

If we take $\xredone$ to be $\beta$-reduction, we easily see that a
consequence of the Church--Rosser property is that if a term has a
normal form, then it is unique. For suppose $M$ can be reduced to $P$
and $Q$, both of which are normal forms. By the Church--Rosser property,
there exists some $N$ such that both $P$ and $Q$ reduce to it. Since
by assumption $P$ and $Q$ are normal forms, the reduction of $P$ and
$Q$ to $N$ can only be the trivial reduction, i.e., $P$, $Q$, and $N$
are identical. This justifies our speaking of \emph{the} normal form
of a term.

In viewing the lambda calculus as a model of computation, then, the
normal form of a term can be thought of as the ``final result'' of the
computation starting with that term. The above corollary means there's
only one, if any, final result of a computation, just like there is
only one result of computing $4 \times (1+2)+3$, namely~$15$.

\begin{thm} \ollabel{thm:str}
  If a relation $\xredone$ satisfies the Church--Rosser property, and $\xred$ is the
  smallest transitive relation containing $\xredone$, then $\xred$ satisfies
  the Church--Rosser property too.
\end{thm}

\begin{proof}
  Suppose 
  \begin{align*}
    M & \xredone P_1 \xredone \dots \xredone P_m \text{ and}\\
    M & \xredone Q_1 \xredone \dots \xredone Q_n.
  \end{align*}
  We will prove the theorem by constructing a grid $N$ of terms of
  height is $m + 1$ and width $n + 1$. We use $N_{i,j}$ to denote the
  term in the $i$-th row and $j$-th column.
  
  We construct $N$ in such a way that $N_{i,j} \xredone N_{i+1,j}$ and
  $N_{i,j} \xredone N_{i,j+1}$. It is defined as follows:
  \begin{align*}
    N_{0,0} &= M \\
    N_{i,0} &= P_i && \text{if } 1 \le i \le m \\
    N_{0,j} &= Q_j && \text{if } 1 \le j \le n \\
  \intertext{and otherwise:}
    N_{i,j} &= R
  \end{align*}
  where $R$ is a term such that $N_{i-1,j} \xredone R$ and $N_{i,j-1}
  \xredone R$. By the Church--Rosser property of $\xredone$, such a
  term always exists.

  Now we have $N_{m,0} \xredone \dots \xredone N_{m,n}$ and $N_{0,n}
  \xredone \dots \xredone N_{m,n}$. Note $N_{m,0}$ is $P$ and $N_{0,n}$
  is $Q$. By definition of $\xred$ the theorem follows.
\end{proof}


% ==== OLP file: content/lambda-calculus/church-rosser/parallel-beta-eta-reduction.tex ====
% Part: lambda-calculus
% Chapter: church-rosser
% Section: parallel-beta-eta-reduction



\olfileid{lam}{cr}{pbe}

\olsection{Parallel $\beta\eta$-reduction}

In this section we prove the Church-Rosser property for parallel
$\beta\eta$-reduction, the parallel reduction notion corresponding to
$\beta\eta$-reduction. 

\begin{defn}[Parallel $\beta\eta$-reduction, $\beredpar$] \ollabel{defn:beredpar}
  \emph{Parallel $\beta\eta$-reduction} ($\beredpar$) on terms is
  inductively defined as follows:
  \begin{enumerate}
    \item \ollabel{defn:beredpar1} $x \beredpar x$.
    \item \ollabel{defn:beredpar2} If $N \xrightarrow{\beta} N'$ then $\lambd[x][N] \beredpar
      \lambd[x][N']$.
    \item \ollabel{defn:beredpar3} If $P \beredpar P'$ and $Q \beredpar Q'$ then $PQ \beredpar
      P'Q'$.
    \item \ollabel{defn:beredpar4} If $N \beredpar N'$ and $Q \beredpar Q'$ then
      $(\lambd[x][N])Q \beredpar \Subst{N'}{Q'}{x}$.
    \item \ollabel{defn:beredpar5} If $N \beredpar N'$ then $\lambd[x][Nx]
      \beredpar N'$, provided $x \notin FV(N)$.
  \end{enumerate}
\end{defn}

\begin{thm}\ollabel{thm:refl}
  $M \beredpar M$.
\end{thm}

\begin{proof}
  Exercise.
\end{proof}

\begin{prob}
  Prove \olref[lam][cr][pbe]{thm:refl}.
\end{prob}

\begin{defn}[$\beta\eta$-complete development]\ollabel{defn:becd}
  The \emph{$\beta\eta$-complete development}~$\becd{M}$ of~$M$ is defined
  as follows:
  \begin{align}
    \becd{x} &= x \ollabel{defn:becd1} \\
    \becd{(\lambd[x][N])} &= \lambd[x][\becd{N}] \ollabel{defn:becd2}\\
    \becd{(PQ)} &= \becd{P}\becd{Q} && \text{if $P$ is not a $\lambd$-abstract} 
    \ollabel{defn:becd3} \\
    \becd{((\lambd[x][N])Q)} &= \Subst{\becd{N}}{\becd{Q}}{x}
                             \ollabel{defn:becd4} \\
    \becd{(\lambd[x][Nx])} &= \becd{N} \ollabel{defn:becd5} & \text{if $x
                                                      \notin FV(N)$}
  \end{align}
\end{defn}

\begin{lem}\ollabel{lem:comp}
  If $M \beredpar M'$ and $R \beredpar R'$, then $\Subst{M}{R}{y}
  \beredpar \Subst{M'}{R'}{y}$.
\end{lem}

\begin{proof}
  By induction on the !!{derivation} of $M \beredpar M'$.

  The first four cases are exactly like those in \olref[pb]{lem:comp}.
  If the last rule is \olref{defn:beredpar5}, then $M$ is
  $\lambd[x][Nx]$, $M'$ is $N'$ for some $x$ and~$N'$ where $x \notin
  FV(N)$, and $N \beredpar N'$. We want to show that
  $\Subst{(\lambd[x][Nx])}{R}{y} \beredpar \Subst{N'}{R'}{y}$, i.e.,
  $\lambd[x][\Subst{N}{R}{y} x] \beredpar \Subst{N'}{R'}{y}$. It
  follows by \olref{defn:beredpar}\olref{defn:beredpar5} and the
  induction hypothesis.
\end{proof}

\begin{lem}\ollabel{lem:cont}
  If $M \beredpar M'$ then $M' \beredpar \becd{M}$.
\end{lem}

\begin{proof}
  By induction on the !!{derivation} of $M \beredpar M'$.

  The first four cases are like those in \olref[pb]{lem:cont}. If the
  last rule is \olref{defn:beredpar5}, then $M$ is $\lambd[x][Nx]$ and
  $M'$ is $N'$ for some $x$, $N$, $N'$ where $x \notin FV(N)$ and $N
  \beredpar N'$. We want to show that $N' \beredpar
  \becd{(\lambd[x][Nx])}$, i.e., $N' \beredpar \becd{N}$, which is
  immediate by induction hypothesis.
\end{proof}

\begin{thm}\ollabel{thm:cr}
  $\beredpar$ has the Church-Rosser property.
\end{thm}

\begin{proof}
  Immediate from \olref{lem:cont}.
\end{proof}

% ==== OLP file: content/lambda-calculus/church-rosser/parallel-beta-reduction.tex ====
% Part: lambda-calculus
% Chapter: church-rosser
% Section: parallel-beta-reduction



\olfileid{lam}{cr}{pb}

\olsection{Parallel $\beta$-reduction}

We introduce the notion of \emph{parallel $\beta$-reduction}, and
prove the it has the Church--Rosser property.

\begin{defn}[parallel $\beta$-reduction, $\bredpar$] \ollabel{defn:bredpar}
  Parallel reduction ($\bredpar$) of terms is inductively defined as follows:
  \begin{enumerate}
    \item \ollabel{defn:bredpar1} $x \bredpar x$.
    \item \ollabel{defn:bredpar2} If $N \xrightarrow{\beta} N'$ then $\lambd[x][N] \bredpar
      \lambd[x][N']$.
    \item \ollabel{defn:bredpar3} If $P \bredpar P'$ and $Q \bredpar Q'$ then $PQ \bredpar
      P'Q'$.
    \item \ollabel{defn:bredpar4} If $N \bredpar N'$ and $Q \bredpar Q'$ then
      $(\lambd[x][N])Q \bredpar \Subst{N'}{Q'}{x}$.
  \end{enumerate}
\end{defn}

Parallel $\beta$-reduction allows us to reduce any number of redices in a
term in one step. It is different from $\beta$-reduction in the sense that we
can only contract redices that occur in the original term, but not redices
arising from parallel $\beta$-reduction. For example, the term
$(\lambd[f][fx])(\lambd[y][y])$ can only be parallel $\beta$-reduced
to itself or to $(\lambd[y][y])x$, but not further to~$x$, although it
$\beta$-reduces to~$x$, because this redex arises only after one step 
of parallel $\beta$-reduction. A second parallel $\beta$-reduction
step yields~$x$, though.

\begin{thm}\ollabel{thm:refl}
  $M \bredpar M$.
\end{thm}

\begin{proof}
  Exercise.
\end{proof}

\begin{prob}
  Prove \olref[lam][cr][pb]{thm:refl}.
\end{prob}

\begin{defn}[$\beta$-complete development]\ollabel{defn:bcd}
  The \emph{$\beta$-complete development}~$\bcd{M}$ of $M$ is defined
  inductively as follows:
  \begin{align}
    \bcd{x} &= x \ollabel{defn:bcd1} \\
    \bcd{(\lambd[x][N])} &= \lambd[x][\bcd{N}] \ollabel{defn:bcd2}\\
    \bcd{(PQ)} &= \bcd{P}\bcd{Q} && \text{if $P$ is not a $\lambd$-abstract} 
    \ollabel{defn:bcd3} \\
    \bcd{((\lambd[x][N])Q)} &= \Subst{\bcd{N}}{\bcd{Q}}{x} \ollabel{defn:bcd4}
  \end{align}
\end{defn}

The $\beta$-complete development of a term, as its name suggests, is a
``complete parallel reduction.'' While for parallel $\beta$-reduction we still
can choose to not contract a redex, for complete development we have
no choice but to contract all of them. Thus the complete development
of $(\lambd[f][fx])(\lambd[y][y])$ is $(\lambd[y][y])x$, not itself.

\begin{editorial}
  This definition has the problem that we haven't introduced how to
  define functions on ($\lambd$-)terms recursively. Will fix in future.
\end{editorial}

\begin{lem}\ollabel{lem:comp}
  If $M \bredpar M'$ and $R \bredpar R'$, then $\Subst{M}{R}{y}
  \bredpar \Subst{M'}{R'}{y}$.
\end{lem}

\begin{proof}
  By induction on the !!{derivation} of $M \bredpar M'$.
  \begin{enumerate}
    \item The last step is \olref{defn:bredpar1}: Exercise.
    \item The last step is \olref{defn:bredpar2}: Then $M$ is 
      $\lambd[x][N]$ and $M'$ is $\lambd[x][N']$,
      where $N \bredpar N'$. We want to prove that
      $\Subst{(\lambd[x][N])}{R}{y} \bredpar
      \Subst{(\lambd[x][N'])}{R'}{y}$, i.e.,
      $\lambd[x][\Subst{N}{R}{y}] \bredpar
      \lambd[x][\Subst{N'}{R}{y}]$. This follows immediately by
      \olref{defn:bredpar2} and the induction hypothesis.
    \item The last step is \olref{defn:bredpar3}: Exercise.
    \item The last step is \olref{defn:bredpar4}: $M$ is
      $(\lambd[x][N])Q$ and $M'$ is $\Subst{N'}{Q'}{x}$. We want to
      prove that $\Subst{((\lambd[x][N])Q)}{R}{y} \bredpar
      \Subst{\Subst{N'}{Q'}{x}}{R'}{y}$, i.e.,
      $(\lambd[x][\Subst{N}{R}{y}])\Subst{Q}{R}{y} \bredpar
      \Subst{\Subst{N'}{R'}{y}}{\Subst{Q'}{R'}{y}}{x}$. This follows
      by \olref{defn:bredpar4} and the induction hypothesis.
  \end{enumerate}
\end{proof}

\begin{prob}
  Complete the proof of \olref[lam][cr][pb]{lem:comp}.
\end{prob}

\begin{lem}\ollabel{lem:cont}
  If $M \bredpar M'$ then $M' \bredpar \bcd{M}$.
\end{lem}
\begin{proof}
  By induction on the !!{derivation} of $M \bredpar M'$.
  \begin{enumerate}
    \item The last rule is \olref{defn:bredpar1}: Exercise.
    \item The last rule is \olref{defn:bredpar2}: $M$ is 
      $\lambd[x][N]$ and $M'$ is $\lambd[x][N']$ with
      $N \bredpar N'$. We want to show that $\lambd[x][N'] \bredpar
      \bcd{(\lambd[x][N])}$, i.e., $\lambd[x][N'] \bredpar
      \lambd[x][\bcd{N}]$ by \olref{defn:bcd2}. It follows by
      \olref{defn:bredpar2} and the induction hypothesis.
    \item The last rule is \olref{defn:bredpar3}:$M$ is $PQ$ and $M'$
      is $P'Q'$ for some $P$, $Q$, $P'$ and $Q'$, with $P \bredpar P'$
      and $Q \bredpar Q'$. By induction hypothesis, we have $P'
      \bredpar \bcd{P}$ and $Q' \bredpar \bcd{Q}$.
      \begin{enumerate}
        \item If $P$ is $\lambd[x][N]$ for some $x$ and $N$, then
          $P'$ must be $\lambd[x][N']$ for some $N'$ with 
          $N \bredpar N'$. By induction hypothesis we have $N' \bredpar \bcd{N}$ and
          $Q' \bredpar \bcd{Q}$. Then $(\lambd[x][N'])Q' \bredpar
          \Subst{\bcd{N}}{\bcd{Q}}{x}$ by \olref{defn:bredpar4}.
        \item If $P$ is not a $\lambd$-abstract, then $P'Q' \bredpar
          \bcd{P}\bcd{Q}$ by \olref{defn:bredpar3}, and the right-hand
          side is $\bcd{PQ}$ by \olref{defn:bcd3}.
      \end{enumerate}
    \item The last rule is \olref{defn:bredpar4}: $M$ is
      $(\lambd[x][N])Q$  and $M'$ is $\Subst{N'}{Q'}{x}$ for some $x$,
      $N$, $Q$, $N'$, and~$Q'$, with $N \bredpar N'$ and $Q \bredpar
      Q'$. By induction hypothesis we know $N' \bredpar \bcd{N}$ and
      $Q' \bredpar \bcd{Q}$. By \olref{lem:comp} we have
      $\Subst{N'}{Q'}{x} \bredpar \Subst{\bcd{N}}{\bcd{Q}}{x}$, the
      right-hand side of which is exactly $\bcd{((\lambd[x][N])Q)}$.
  \end{enumerate}
\end{proof}

\begin{prob}
  Complete the proof of \olref[lam][cr][pb]{lem:cont}.
\end{prob}

\begin{thm}\ollabel{thm:cr}
  $\bredpar$ has the Church--Rosser property.
\end{thm}

\begin{proof}
  Immediate from \olref{lem:cont}.
\end{proof}



